\lecture{9.4}{Sensitivity Analysis for Unmeasured Confounding}
\label{lec:9-4}

\subsection{Violation of Conditional Ignorability}

For a binary treatment, the usual identification proof for $\E\{\Yone\}$ starts
from
\[
  \E\{\Yone\}
  =
  \sum_w\sum_a
  \E(\Yone\given A=a,W=w)\Prob(A=a\given W=w)\Prob(W=w).
\]
Split the sum over $a$:
\[
\begin{aligned}
  \E\{\Yone\}
  &=
  \sum_w \E(\Yone\given A=1,W=w)\Prob(A=1\given W=w)\Prob(W=w)\\
  &\quad+
  \sum_w \E(\Yone\given A=0,W=w)\Prob(A=0\given W=w)\Prob(W=w).
\end{aligned}
\]
The first conditional mean is identified by consistency:
\[
  \E(\Yone\given A=1,W=w)=\E(Y\given A=1,W=w).
\]
The second conditional mean concerns untreated individuals under treatment:
\[
  \E(\Yone\given A=0,W=w).
\]
It is not observed directly.

If $\Yone\ci A\given W$, then
\[
  \E(\Yone\given A=0,W=w)
  =
  \E(\Yone\given A=1,W=w)
  =
  \E(Y\given A=1,W=w),
\]
and the usual adjustment formula follows. If conditional ignorability fails,
this replacement is not justified.

\subsection{Non-Identifiable Counterfactual Components}

The unidentified piece is the distribution of $\Yone$ among those with $A=0$.
Similarly, identification of $\E\{\Yzero\}$ requires information about $\Yzero$
among those with $A=1$. These are counterfactual components that the observed
data do not reveal without assumptions.

\begin{keyidea}
Unmeasured confounding breaks identification because some counterfactual
outcome distributions are never observed in the treatment groups where they are
needed.
\end{keyidea}

\subsection{Bounds Without Sensitivity Parameters}

If $Y$ is bounded, partial identification is possible without specifying a
sensitivity model. Suppose $Y\in[0,1]$. Then
\[
  0\le \E(\Yone\given A=0,W=w)\le 1.
\]
Therefore $\E\{\Yone\}$ is bounded by replacing the unidentified component with
0 or 1:
\[
\begin{aligned}
  \sum_w &\E(Y\given A=1,W=w)\Prob(A=1\given W=w)\Prob(W=w)\\
  &\le \E\{\Yone\}\\
  &\le
  \sum_w
  \left[
  \E(Y\given A=1,W=w)\Prob(A=1\given W=w)
  +\Prob(A=0\given W=w)
  \right]\Prob(W=w).
\end{aligned}
\]
These bounds are honest but often wide.

\subsection{Bounds With Sensitivity Parameters}

Sensitivity parameters narrow the set of plausible counterfactual distributions
by restricting the magnitude of unmeasured confounding. For example, one might
bound the difference
\[
  \left|
  \E(\Yone\given A=0,W=w)-\E(\Yone\given A=1,W=w)
  \right|
  \le \delta.
\]
For each value of $\delta$, this yields an interval of possible values for
$\E\{\Yone\}$ and therefore for the ATE. Reporting the range over plausible
$\delta$ values makes the dependence on untestable assumptions explicit.

\subsection{Point Identification With Sensitivity Parameters}

Another approach specifies a sensitivity model that maps observed outcome
distributions to the unobserved counterfactual distributions. Once the
sensitivity parameter is fixed, the causal parameter becomes identified.

\subsection{Robins Sensitivity Model}

One useful model relates the density of $\Yone$ among untreated individuals to
the density among treated individuals:
\[
  p(y_1\given A=0,W)
  =
  \frac{
  p(y_1\given A=1,W)\exp\{\gamma S(y_1)\}
  }{
  \E[\exp\{\gamma S(Y)\}\given A=1,W]
  }.
\]
Here $S(\cdot)$ is a specified function and $\gamma$ is a sensitivity parameter.
When $\gamma=0$, the two conditional distributions are equal, corresponding to
no unmeasured confounding for $\Yone$ within strata of $W$.

Under this model,
\[
  \E(\Yone\given A=0,W=w)
  =
  \frac{
  \E[Y\exp\{\gamma S(Y)\}\given A=1,W=w]
  }{
  \E[\exp\{\gamma S(Y)\}\given A=1,W=w]
  }.
\]
Thus, for fixed $\gamma$,
\[
\begin{aligned}
  \E\{\Yone\}
  &=
  \sum_w \E(Y\given A=1,W=w)\Prob(A=1\given W=w)\Prob(W=w)\\
  &\quad+
  \sum_w
  \frac{
  \E[Y\exp\{\gamma S(Y)\}\given A=1,W=w]
  }{
  \E[\exp\{\gamma S(Y)\}\given A=1,W=w]
  }
  \Prob(A=0\given W=w)\Prob(W=w).
\end{aligned}
\]

\begin{estimationbox}
A sensitivity analysis workflow is:
\begin{enumerate}
  \item Choose a sensitivity model and interpret its parameter scale.
  \item Select a scientifically plausible grid of sensitivity values.
  \item Estimate the identified functional at each value.
  \item Report how the causal conclusion changes across the grid.
\end{enumerate}
\end{estimationbox}
